\DocumentMetadata{
  pdfversion=2.0,
  pdfstandard=ua-2,
  lang=en-US,
  tagging=on,
  tagging-setup={math/setup=mathml-SE,tabsorder=structure}
}
\documentclass[11pt,letterpaper]{article}
\usepackage[margin=0.68in,includefoot]{geometry}
\usepackage{newcomputermodern}
\usepackage{amsmath,amscd,mathtools}
\providecommand{\square}{\mdlgwhtsquare}
\usepackage[hidelinks]{hyperref}
\hypersetup{
  pdftitle={Partial Differential Equations PhD Exam, May 2001},
  pdfauthor={University of Florida Department of Mathematics},
  pdfsubject={Graduate mathematics examination},
  pdfdisplaydoctitle=true,
  bookmarksopen=true
}
\setcounter{secnumdepth}{0}
\setlength{\parindent}{0pt}
\setlength{\parskip}{0.28em}
\setlength{\emergencystretch}{3em}
\setlength{\leftmargini}{2.15em}
\setlength{\leftmarginii}{2.65em}
\setlength{\leftmarginiii}{3em}
\renewcommand{\labelenumi}{\textbf{\theenumi.}}
\pagestyle{empty}
\raggedbottom
\allowdisplaybreaks
\newenvironment{examproblems}[1]{%
  \begin{enumerate}%
  \setcounter{enumi}{#1}%
  \setlength{\topsep}{0.35em}%
  \setlength{\partopsep}{0pt}%
  \setlength{\itemsep}{0.48em}%
  \setlength{\parsep}{0.18em}%
}{\end{enumerate}}
\newenvironment{examsubparts}{%
  \begin{enumerate}%
  \setlength{\topsep}{0.18em}%
  \setlength{\partopsep}{0pt}%
  \setlength{\itemsep}{0.16em}%
  \setlength{\parsep}{0.1em}%
}{\end{enumerate}}
\newenvironment{examinstructions}{%
  \par\begingroup\itshape
}{%
  \par\endgroup\vspace{0.18em}
}
\makeatletter
\renewcommand\section{\@startsection{section}{1}{0pt}%
  {-1.25ex plus -.25ex minus -.1ex}%
  {0.55ex plus .1ex}%
  {\normalfont\large\bfseries}}
\renewcommand{\@maketitle}{%
  \newpage\null\vskip -1.2em%
  {\footnotesize\bfseries
    \href{https://www.ufl.edu/}{University of Florida}, \href{https://math.ufl.edu/}{Department of Mathematics}\par}%
  \vskip 0.18em%
  \tagstructbegin{tag=Title}%
    {\LARGE\bfseries\href{https://gma.math.ufl.edu/exams/pde-phd/}{Partial Differential Equations PhD Exam}, May 2001\par}%
  \tagstructend%
  \vskip 0.35em%
  \tagmcbegin{artifact}%
    \hrule height 1pt%
  \tagmcend%
  \vskip 0.55em%
}
\makeatother
\title{Partial Differential Equations PhD Exam, May 2001}
\author{}
\date{}
\begin{document}
\maketitle
\thispagestyle{empty}
\begin{examinstructions}
Do all problems 1–4. Choose one problem of 5 and 6 and one problem of 7 and 8.
\end{examinstructions}
\begin{examproblems}{0}
\item
Let~\(B\) be the unit ball in \(\mathbb{R}^3\). For which values of~\(\alpha\) does the function \(|x|^\alpha\) belong to \(H^1(B)\)? Justify your answer.
\item
Let \(U\subset\mathbb{R}^n\) be open and bounded. Let \(u_m\) and \(v_m\) be bounded sequences in \(H^1(U)\). Show that there exist subsequences \(u_{m_j}\) of \(u_m\) and \(v_{m_j}\) of \(v_m\), and functions \(u,v\in H^1(U)\), such that 
\[
\int_U u_{m_j}v_{m_j}\,dx\to\int_U uv\,dx,
\]
 and 
\[
\int_U D_{x_i}(u_{m_j}v_{m_j})\,dx\to\int_U D_{x_i}(uv)\,dx,\qquad i=1,\ldots,n,
\]
 as \(j\to\infty\).
\item
Let \(U\subset\mathbb{R}^n\) be open, bounded, and connected with smooth boundary \(\partial U\). Let \(u(x,t)\) be a smooth solution to the IBV problem: 
\[
u_t=\Delta u,\qquad (x,t)\in U_T,
\]
 
\[
\frac{\partial u}{\partial n}=0,\qquad (x,t)\in\partial U\times[0,T],
\]
 
\[
u(x,0)=f(x),\qquad x\in U.
\]
 Let 
\[
(u)_U(t)=\frac{1}{|U|}\int_U u(x,t)\,dx.
\]
 Show that the following statements hold.

\par
Hint: Consider the equation for \(v=u-(u)_U\), and use energy estimate for the equation of~\(v\), and Poincaré’s inequality.
\begin{examsubparts}
\item[\textnormal{(a)}]
\(\displaystyle \frac{d}{dt}(u)_U(t)\equiv0.\)
\item[\textnormal{(b)}]
\(u\to(u)_U\) in \(L^2(U)\), as \(t\to\infty\).
\end{examsubparts}
\item
Let \(U\subset\mathbb{R}^n\) be open, bounded, and connected with smooth boundary \(\partial U\). Let~\(L\) be a uniformly elliptic differential operator of the form 
\[
Lu=-\sum_{i,j=1}^n a_{ij}u_{x_i x_j}.
\]
 Suppose that~\(f\) is a bounded function and \(u,v\in C^2(\overline U)\) satisfy 
\[
Lu=f,\qquad Lv\geq1,\qquad x\in U,
\]
 
\[
u(x)=0,\qquad v(x)\geq0,\qquad x\in\partial U.
\]
 Suppose that there exists a point \(x_0\in\partial U\) such that \(v(x_0)=0\). Prove that there exists a constant \(C>0\) such that 
\[
|Du(x_0)|\leq C\left|\frac{\partial v}{\partial n}(x_0)\right|.
\]
\item
Let \(U\subset\mathbb{R}^n\) be open and bounded with smooth boundary \(\partial U\). We say that a function \(u\in H_0^2(U)\) is a weak solution of the biharmonic equation 
\[
\Delta^2u=f,\qquad x\in U, \tag{5.1}
\]
 
\[
u=0,\qquad \frac{\partial u}{\partial n}=0,\qquad x\in\partial U, \tag{5.2}
\]
 if 
\[
\int_U\Delta u\,\Delta v\,dx=\int_U fv\,dx,\qquad\text{for any }v\in H_0^2(U). \tag{5.3}
\]
 Prove the following statements:
\begin{examsubparts}
\item[\textnormal{(a)}]
For a given \(f\in L^2(U)\) there is an unique solution to the problem (5.3).
\item[\textnormal{(b)}]
If \(u\in C^4(\overline U)\) and \(f\in C^0(\overline U)\) satisfy (5.3), then~\(u\) satisfies (5.1) at each point \(x\in U\).
\end{examsubparts}
\item
Let \(U\subset\mathbb{R}^n\) be open and bounded. Show that \(u\in H_0^1(U)\) is a weak solution to the boundary value problem 
\[
-\Delta u+u=1,\qquad x\in U,
\]
 
\[
u=0,\qquad x\in\partial U,
\]
 if and only if~\(u\) minimizes 
\[
E(v)=\int_U\bigl(|\nabla v|^2+v^2-2v\bigr)\,dx
\]
 over \(v\in H_0^1(U)\).
\item
Let \(U\subset\mathbb{R}^n\) be open and bounded. We say \(v\in C^2(\overline U)\) is subharmonic if \(-\Delta v\leq0\) for all \(x\in U\).
\begin{examsubparts}
\item[\textnormal{(a)}]
Prove that for any subharmonic~\(v\), 
\[
v(x)\leq\frac{1}{|B(x,r)|}\int_{B(x,r)}v(y)\,dy,\qquad\text{for all }B(x,r)\subset U.
\]
\item[\textnormal{(b)}]
Let \(\phi:\mathbb{R}\to\mathbb{R}\) be smooth and convex. Suppose that~\(u\) is harmonic and \(v=\phi(u)\). Prove that~\(v\) is subharmonic.
\end{examsubparts}
\item
Let \(U\subset\mathbb{R}^n\) be open and bounded with smooth boundary \(\partial U\).
\begin{examsubparts}
\item[\textnormal{(a)}]
Write the fundamental solution \(\Phi(x)\) of Laplace’s equation.
\item[\textnormal{(b)}]
Prove that for any point \(x\in U\) and any function \(u\in C^2(\overline U)\), 
\[
u(x)=\int_U\Phi(y-x)\Delta u(y)\,dy
+\int_{\partial U}\left[\Phi(y-x)\frac{\partial u}{\partial n}(y)-u(y)\frac{\partial\Phi}{\partial n}(y-x)\right]dS(y),
\]
 where~\(\Phi\) is the fundamental solution of Laplace’s equation, and~\(n\) is the outward unit normal to \(\partial U\).
\end{examsubparts}
\end{examproblems}
\end{document}
